\chapter{Le marché du Trésor}\label{ch:m2:the-treasury-market}

À 9:33:19, heure de New York, le 15 octobre 2014, le rendement de la note du
Trésor américain à dix ans s'établissait à 2.02~\%. Six minutes plus tard, à 9:39:39, il était à 1.86~\%~;
à 9:44:35, il remontait à 1.99~\%. Aucune publication économique, aucune déclaration,
ni aucune dépêche n'a expliqué l'un ou l'autre mouvement. La note sert de référence pour la
courbe des rendements en dollars, et pendant douze minutes elle s'est comportée comme une petite action. Le rapport officiel qui a suivi a établi que plus
de la moitié des transactions de la journée, dans les carnets d'ordres électroniques du marché au comptant
et sur les contrats à terme, avaient été réalisées par des sociétés de trading pour compte propre, et non
par les banques qui avaient autrefois fait le marché. Ce chapitre décrit comment
les titres du Trésor sont vendus, où ils se négocient, qui les négocient, et ce que deux des
journées les plus étudiées du marché ont révélé sur sa structure.

\section{Adjudications et spécialistes en valeurs du Trésor}

Le Trésor vend ses titres lors d'adjudications, selon un calendrier publié. Les
soumissions portent sur des rendements, et non sur des prix~; le coupon d'une nouvelle note n'est fixé qu'ensuite.

\begin{definition}[Adjudication à prix uniforme]\label{def:m2:the-treasury-market:auction}
Dans une \emph{adjudication à prix uniforme}\index{adjudication à prix uniforme} (uniform-price auction, ou adjudication à prix unique)
les soumissions non concurrentielles, pour un montant limité, sont d'abord servies
au rendement qui en résulte~; les soumissions concurrentielles, chacune
composée d'un rendement et d'un montant, sont ensuite retenues, du plus bas rendement vers le plus haut, jusqu'à épuisement de
l'émission. Le rendement le plus élevé retenu est le \emph{rendement de coupure} (stop-out)~:
tous les soumissionnaires retenus le paient, et les soumissions situées exactement à ce niveau sont servies au
prorata.
\end{definition}

\begin{method}[Déroulement d'une adjudication]\label{met:m2:the-treasury-market:run}
(1) Plafonner les soumissions reconnues de chaque soumissionnaire à 35~\% de l'émission, en commençant par les
rendements les plus bas. (2) Soustraire de l'émission les attributions non concurrentielles. (3) Trier
les soumissions concurrentielles par rendement et en cumuler les montants~; le rendement auquel le montant cumulé
atteint le solde est le rendement de coupure. (4) Servir intégralement les soumissions situées sous ce niveau~;
servir les soumissions situées à ce niveau au prorata du pourcentage qui complète l'émission, arrondi au
centième de pour cent. (5) Fixer le coupon au huitième de pour cent qui
donne un prix le plus proche possible, sans le dépasser, du pair au rendement de coupure.
\end{method}

\begin{definition}[Spécialiste en valeurs du Trésor]\label{def:m2:the-treasury-market:pd}
Un \emph{spécialiste en valeurs du Trésor}\index{spécialiste en valeurs du Trésor} (primary dealer) est une société que la
Federal Reserve Bank of New York a désignée comme contrepartie de ses opérations,
en échange d'engagements~: soumissionner au prorata dans chaque adjudication du Trésor
à des prix raisonnablement compétitifs, assurer la tenue de marché pour la Banque
pour le compte de ses titulaires de comptes officiels, et rendre compte de son activité.
\end{definition}

L'obligation de soumissionner fait des spécialistes les preneurs fermes de l'adjudication~: si
les investisseurs restent à l'écart, les soumissions des spécialistes remplissent l'émission, à un rendement
assez élevé pour valoir leur risque. La part d'une adjudication que les spécialistes ont dû absorber se lit donc
comme les autres statistiques de l'adjudication, en tant que mesure de la demande.

\begin{definition}[Taux de couverture et tail d'adjudication]\label{def:m2:the-treasury-market:tail}
Le \emph{taux de couverture}\index{taux de couverture} (bid-to-cover) est le montant soumis
divisé par le montant offert. Le \emph{tail d'adjudication}\index{tail d'adjudication} (auction tail)
est le rendement de coupure moins le rendement du \omterm{def:m2:the-treasury-market:wi}{marché gris} (when-issued)
(\cref{def:m2:the-treasury-market:wi}) à l'heure limite de soumission~: si l'écart est positif,
l'adjudication s'est conclue à un rendement supérieur à celui du marché (elle a \emph{tailé})~;
si l'écart est négatif, elle s'est conclue à un rendement inférieur (elle a \emph{stoppé through}).
\end{definition}

\begin{example}[Une adjudication simulée à dix ans]\label{ex:m2:the-treasury-market:auction}
Sont offerts 42~milliards de dollars de notes à dix ans~; 300~millions de dollars de
soumissions non concurrentielles laissent 41.7~milliards de dollars pour 76.1~milliards de dollars de soumissions
concurrentielles. La note en \omterm{def:m2:the-treasury-market:wi}{marché gris} se traite à 4.180~\% à 13:00.
Les soumissions de la \cref{fig:m2:the-treasury-market:demand} atteignent 41.7~milliards de dollars à 4.198~\%~:
c'est le rendement de coupure, les soumissions à ce niveau sont servies à 95.02~\%, le tail est de 1.8 points de base
et le \omterm{def:m2:the-treasury-market:tail}{taux de couverture} de 1.82. Les soumissionnaires indirects (investisseurs passant par un
spécialiste) en prennent 51.1~\%, les soumissionnaires directs 10.6~\% et les spécialistes 38.3~\%. Le coupon est
fixé à 4.125~\%, et chaque gagnant paie 99.409 pour une note que le marché valorisait
à 4.180~\% une minute plus tôt. L'adjudication est simulée~; la procédure est celle
du Trésor.
\end{example}

\begin{omfigure}
\begin{tikzpicture}
\begin{axis}[width=11cm,height=5cm,
  xlabel={rendement soumis (\%)},ylabel={soumissions cumulées (Md USD)},
  tick label style={font=\small},label style={font=\small},
  xmin=4.17,xmax=4.225,ymin=0,ymax=80,grid=major,grid style={gray!25},
  xtick={4.17,4.18,4.19,4.2,4.21,4.22},
  xticklabel style={/pgf/number format/fixed,/pgf/number format/precision=2},
  table/col sep=comma]
\addplot[omDef,very thick,const plot] table[x=yield,y=cumbn]{figdata/markets-2/04-the-treasury-market/demand.csv};
\draw[omThm,thick,dashed] (axis cs:4.17,41.7) -- (axis cs:4.225,41.7);
\draw[omProp,thick] (axis cs:4.198,0) -- (axis cs:4.198,80);
\draw[gray,thick,densely dotted] (axis cs:4.180,0) -- (axis cs:4.180,80);
  \node[font=\footnotesize,anchor=north east,text=omThm] at (axis cs:4.2245,41) {à vendre~: 41.7};
  \node[font=\footnotesize,anchor=north west,text=omProp] at (axis cs:4.1985,76) {rendement de coupure 4.198};
  \node[font=\footnotesize,anchor=north west,text=gray] at (axis cs:4.1805,76) {marché gris 4.180};
\end{axis}
\end{tikzpicture}

\omcaption{La courbe de demande de l'\cref{ex:m2:the-treasury-market:auction}~: soumissions
cumulées à partir du rendement le plus bas. Le rendement de coupure est là où la courbe coupe le
montant à vendre~; l'écart entre ce rendement et le rendement du \omterm{def:m2:the-treasury-market:wi}{marché gris} à la
date limite est le tail. Un carnet dont les soumissions se situent plus près du niveau du \omterm{def:m2:the-treasury-market:wi}{marché gris}
stopperait through. Données~: le carnet de soumissions simulé du chapitre.}\label{fig:m2:the-treasury-market:demand}
\end{omfigure}

\begin{dated}{2026-09}{Adjudications du Trésor et spécialistes en valeurs du Trésor}\label{dat:m2:the-treasury-market:auctions}
Adjudications à prix unique pour tous les titres négociables depuis novembre 1998.
Soumissions non concurrentielles jusqu'à 10~millions de dollars par soumissionnaire, servies intégralement. Aucun
soumissionnaire concurrentiel n'est reconnu pour plus de 35~\% de l'émission. Les soumissions
clôturent habituellement à 12:00 (non concurrentielles) et à 13:00 (concurrentielles), heure de
New York. La Federal Reserve Bank of New York recense 26 \omterm{def:m2:the-treasury-market:pd}{spécialistes en valeurs du Trésor}, dont le
plus récent a été ajouté le 15 janvier 2026.
\end{dated}

\section{Marché gris, on-the-run et off-the-run}

\begin{definition}[Marché gris]\label{def:m2:the-treasury-market:wi}
\emph{marché gris}\index{marché gris} (when-issued trading) désigne la négociation d'un titre
qui a été annoncé mais n'a pas encore été adjugé ni émis, pour un règlement-livraison
à sa date d'émission. Le coupon n'étant pas encore connu, le titre est coté en rendement.
\end{definition}

\begin{definition}[On-the-run et off-the-run]\label{def:m2:the-treasury-market:otr}
Le titre le plus récemment adjugé pour chaque échéance d'origine est
\emph{on-the-run}\index{on-the-run} (l'indice de référence)~; toutes les émissions antérieures de cette
échéance sont \emph{off-the-run}\index{off-the-run}. Chaque nouvelle adjudication fait passer
l'ancien indice de référence en off-the-run.
\end{definition}

Le \omterm{def:m2:the-treasury-market:wi}{marché gris} permet aux dealers de prévendre l'adjudication à leurs clients et de couvrir
leurs soumissions~: un dealer qui s'attend à se voir attribuer un milliard vend à découvert un
milliard en \omterm{def:m2:the-treasury-market:wi}{marché gris} et le rachète lors de l'adjudication. Il rend aussi le \omterm{def:m2:the-treasury-market:tail}{tail d'adjudication} observable,
puisqu'il existe à l'échéance un cours de marché à comparer au
résultat. Une fois émis, le nouveau titre devient la référence du marché~: l'essentiel
des négociations électroniques, des couvertures et de la specialness repo
(\cref{ch:m2:repo-and-specials}) se concentre sur lui. Son rendement est en général un peu
inférieur à celui d'un titre \omterm{def:m2:the-treasury-market:otr}{off-the-run} d'échéance presque identique. L'écart rémunère
la liquidité et le financement moins coûteux qu'apporte sa specialness, et
le négocier, en étant long du \omterm{def:m2:the-treasury-market:otr}{off-the-run} bon marché contre l'indice de référence cher, est l'une des
plus anciennes opérations de valeur relative du marché (livre~9).

\section{Courtiers interdealers et plateformes dealer-to-client}

\begin{definition}[Courtier interdealers et plateforme dealer-to-client]\label{def:m2:the-treasury-market:idb}
Un \emph{courtier interdealers}\index{courtier interdealers} (inter-dealer broker) exploite la plateforme de négociation sur
lesquelles les dealers, et de plus en plus des non-dealers, négocient entre eux,
de manière anonyme~: pour les \omterm{def:m2:money-markets:bill}{bons du Trésor} de référence, un carnet d'ordres à cours limité central~;
pour les émissions \omterm{def:m2:the-treasury-market:otr}{off-the-run}, surtout des courtiers assistés par la voix. Une \emph{plateforme
dealer-to-client}\index{plateforme dealer-to-client} permet à un client de demander des cotations à
plusieurs dealers à la fois et de négocier au meilleur prix (\cref{ch:m2:how-bonds-trade}).
\end{definition}

Les deux niveaux ne fixent pas leurs prix de la même façon. Dans les carnets d'ordres, un prix est ferme et
quiconque le frappe négocie~; au niveau client, un dealer cote pour un client qu'il connaît,
pour une taille que le client précise, et couvre le résultat sur le carnet d'ordres.
En 2014, les carnets d'ordres interdealers pour les titres de référence étaient BrokerTec
et eSpeed. Le premier propose un \emph{workup}~: après une transaction dans le carnet, tous les
participants peuvent négocier davantage à ce prix pendant une courte fenêtre, et le rapport
a constaté que la plus grande part du volume de BrokerTec s'est produite dans les workups.

\begin{omfigure}
\begin{tikzpicture}[font=\small,
  box/.style={draw,thick,rounded corners=2pt,align=center,minimum height=0.95cm,inner sep=4pt},
  arr/.style={-{Stealth[length=2.5mm]},thick}]
\node[box,fill=omProp!8,draw=omProp,minimum width=2.6cm] (tsy) at (0,2.4) {Trésor\\(adjudications)};
\node[box,fill=omDef!8,draw=omDef,minimum width=2.8cm] (pd) at (4.6,2.4) {spécialistes en valeurs du Trésor\\et soumissionnaires directs};
\node[box,fill=omThm!8,draw=omThm,minimum width=3.2cm] (idb) at (4.6,0) {carnets d'ordres interdealers\\(dealers et PTF)};
\node[box,fill=gray!8,draw=gray,minimum width=2.8cm] (fut) at (0,0) {bourse de contrats à terme};
\node[box,fill=omDef!8,draw=omDef,minimum width=2.8cm] (cli) at (9.6,2.4) {clients~: fonds,\\banques centrales, banques};
\node[box,fill=omProp!8,draw=omProp,minimum width=2.8cm] (ccp) at (9.6,0) {chambre de compensation};
\draw[arr] (tsy) -- node[above=11pt,font=\footnotesize] {nouvelles émissions} (pd);
\draw[arr,<->] (pd) -- node[right,font=\footnotesize] {couverture, négociation} (idb);
\draw[arr,<->] (pd) -- node[above=11pt,font=\footnotesize,align=center] {cotations sur demande} (cli);
\draw[arr,<->] (fut) -- node[below=11pt,font=\footnotesize] {base} (idb);
\draw[arr,dashed] (idb) -- node[below=11pt,font=\footnotesize] {transactions compensées} (ccp);
\end{tikzpicture}

\omcaption{La structure du marché américain du Trésor. Le Trésor vend aux
spécialistes et soumissionnaires directs lors de l'adjudication~; les clients achètent surtout auprès des dealers
sur demande~; les dealers et les sociétés de trading pour compte propre négocient entre eux dans les
carnets d'ordres interdealers et sur les contrats à terme. L'obligation de compensation
(\cref{dat:m2:the-treasury-market:clearing}) fera passer une plus grande part des négociations secondaires
par une chambre de compensation.}\label{fig:m2:the-treasury-market:structure}
\end{omfigure}

\section{Sociétés de trading pour compte propre et marché électronique}

Les sociétés de trading pour compte propre (\emph{proprietary trading firms}, livre~1,
chapitre 1) négocient leur propre capital, surtout par voie électronique, surtout sur de courts
horizons, et surtout sur les titres de référence et les contrats à terme. Le 15
octobre 2014, elles représentaient plus de la moitié du volume dans le
marché au comptant interdealers et les contrats à terme analysés par les régulateurs, et
leur activité était concentrée~: les dix plus actives ont réalisé plus de 90~\% des
négociations des sociétés de trading pour compte propre sur le marché au comptant. Elles fournissaient l'essentiel
de la profondeur affichée, et la retiraient quand elles le voulaient. Un marché dont la
liquidité est fournie par des sociétés sans obligation de la fournir, l'obligation d'adjudication des
dealers n'ayant pas de contrepartie sur le marché secondaire, est profond en temps normal
et peut devenir mince les jours où on en a besoin.

\section{15 octobre 2014 et mars 2020}

\begin{definition}[Flash rally]\label{def:m2:the-treasury-market:flash}
Un \emph{flash rally}\index{flash rally} est une hausse rapide et importante des cours (baisse
des rendements) qui s'inverse en quelques minutes, sans nouvelle pour expliquer l'un ou l'autre mouvement,
l'image miroir d'un flash crash (livre~1, chapitre~31).
\end{definition}

Les douze minutes du titre à dix ans, le 15 octobre 2014
(\cref{fig:m2:the-treasury-market:oct2014}) ont suivi une publication économique faible sur les ventes au détail à 8:30,
qui avait déjà fait baisser les rendements et retiré la profondeur des carnets. Dans la fenêtre
d'événement, les carnets d'ordres ne représentaient plus qu'un cinquième de leur profondeur habituelle
sur le dix ans, le volume des transactions a bondi, et l'autonégociation de sociétés isolées
était anormalement élevée. Le rapport n'a relevé aucune cause unique. Le rendement a clôturé à
2.14~\%, soit six points de base sous la clôture précédente, après une amplitude de 37
points de base, une amplitude observée seulement trois fois depuis 1998, et chaque fois à la suite
d'une annonce importante.

\begin{omfigure}
\begin{tikzpicture}[font=\small,x=0.9cm,y=0.9cm]
\draw[->,thick] (0,-0.9) -- (12.5,-0.9) node[right,font=\footnotesize] {heure (ET)};
\draw[->,thick] (0,-0.9) -- (0,4.7) node[above,font=\footnotesize] {rendement à dix ans (\%)};
\foreach \y/\l in {0.5/1.85,1.75/1.95,3.0/2.05,4.25/2.15}{\draw (0.1,\y) -- (-0.1,\y) node[left,font=\footnotesize] {\l};}
% y = 0.5 + 12.5 (yield - 1.85); x = 1 + 0.85 (minutes after 9:33:19), so time is to scale
\coordinate (a) at (1.0,2.625);    % 2.02 at 9:33:19
\coordinate (b) at (6.383,0.625);  % 1.86 at 9:39:39 (6 min 20 s later)
\coordinate (c) at (10.578,2.25);  % 1.99 at 9:44:35 (4 min 56 s after that)
\draw[omDef,very thick] (a) -- (b) -- (c);
\foreach \p/\t/\v/\pos in {a/9:33:19/2.02/above,b/9:39:39/1.86/below,c/9:44:35/1.99/above}{
  \fill[omDef] (\p) circle (2.2pt);
  \node[\pos=5pt,font=\footnotesize,align=center] at (\p) {\t\\\v\%};}
\draw[omProp,dashed] (0,4.125) -- (12.3,4.125) node[below left,font=\footnotesize,text=omProp] {clôture 2.14~\%};
\end{tikzpicture}

\omcaption{La fenêtre d'événement du 15 octobre 2014~: les trois rendements et horaires
publiés dans le rapport conjoint des services des régulateurs, reliés par des segments droits (le
trajet entre eux n'était pas lisse), et la clôture de la journée. Seize points de base
en baisse en six minutes, puis remontée, dans l'indice de référence de la
courbe des rendements en dollars. Données~: Joint Staff Report, 2015.}\label{fig:m2:the-treasury-market:oct2014}
\end{omfigure}

\begin{definition}[Ruée vers la liquidité]\label{def:m2:the-treasury-market:dash}
Une \emph{ruée vers la liquidité}\index{ruée vers la liquidité} (dash for cash) est une précipitation d'investisseurs de toute nature
à vendre des actifs liquides, y compris les obligations d'État les plus sûres, pour se procurer des liquidités,
de sorte que les titres habituellement achetés en période de crise sont vendus à la place.
\end{definition}

En mars 2020, alors que la pandémie fermait les économies, des investisseurs de toutes sortes
cherchaient des liquidités. L'examen international de l'épisode a constaté que d'importantes
ventes de \omterm{def:m2:money-markets:bill}{bons du Trésor} américain par certains investisseurs non bancaires à effet de levier et par des détenteurs étrangers
ont rencontré des dealers qui ne pouvaient pas, ou ne voulaient pas, les absorber~; la combinaison
est devenue auto-renforçante, et des prix qui évoluent normalement de concert, les titres \omterm{def:m2:the-treasury-market:otr}{on-the-run}
et \omterm{def:m2:the-treasury-market:otr}{off-the-run}, les \omterm{def:m2:money-markets:bill}{bons du Trésor} et leurs contrats à terme, se sont dissociés. Le 15 mars, la Réserve fédérale a ramené son taux à zéro
et annoncé des achats d'au moins 500~milliards de dollars de \omterm{def:m2:money-markets:bill}{bons du Trésor} américain~; le 23 mars,
elle a remplacé la consigne par des achats «~selon les montants nécessaires~» pour le
bon fonctionnement du marché. Ses avoirs sont passés de 2\,523~milliards de dollars le 11 mars à
3\,789~milliards de dollars le 15 avril, soit 1\,266~milliards en cinq semaines
(\cref{fig:m2:the-treasury-market:fed2020}). La banque centrale est devenue, pendant un
mois, le dealer de dernier recours du marché, un rôle que l'obligation de compensation et
les autres réformes des régulateurs depuis lors visent à le rendre moins nécessaire.

\begin{omfigure}
\begin{tikzpicture}
\begin{axis}[width=11cm,height=4.8cm,
  ylabel={milliards de dollars},
  tick label style={font=\small},label style={font=\small},
  xmin=0,xmax=238,ymin=2000,ymax=4500,grid=major,grid style={gray!25},
  xtick={28,59,88,119,149,180,210},xticklabels={janv.,févr.,mars,avr.,mai,juin,juil.},
  table/col sep=comma]
\addplot[omDef,very thick,mark=*,mark size=1.1pt] table[x=t,y=usdbn]{figdata/markets-2/04-the-treasury-market/fed2020.csv};
\draw[omThm,dashed] (axis cs:102,2000) -- (axis cs:102,4500);
\draw[omThm,dashed] (axis cs:110,2000) -- (axis cs:110,4500);
\node[font=\footnotesize,anchor=east,text=omThm] at (axis cs:101,4250) {15 mars};
\node[font=\footnotesize,anchor=west,text=omThm] at (axis cs:111,4250) {23 mars};
\end{axis}
\end{tikzpicture}

\omcaption{Titres du Trésor détenus par la Réserve fédérale, en données hebdomadaires, de décembre
2019 à juillet 2020. Les deux traits pointillés marquent l'annonce d'achats d'au moins 500
milliards de dollars et le passage à des achats «~selon les montants nécessaires~». Données~:
Federal Reserve H.4.1, via la série FRED TREAST.}\label{fig:m2:the-treasury-market:fed2020}
\end{omfigure}

\begin{dated}{2026-09}{Obligation de compensation des bons du Trésor américain}\label{dat:m2:the-treasury-market:clearing}
Adoptée par la SEC en décembre 2023, la règle impose aux membres des agences de compensation du
Trésor de compenser les transactions éligibles du marché secondaire~: les transactions au comptant
éligibles à partir du 31 décembre 2026 et les opérations de repo éligibles à partir du 30 juin 2027 (dates
prolongées d'un an en février 2025). Le 22 septembre 2026, un commissaire de la SEC
a déclaré que la SEC n'entendait pas les prolonger à nouveau. Deux agences de compensation ont été
approuvées pour agir en tant que contreparties centrales des transactions au comptant et des repos sur les \omterm{def:m2:money-markets:bill}{bons du Trésor}.
\end{dated}

\section{Tutoriel~: conduire une adjudication}\label{tut:m2:the-treasury-market:auction}

\begin{tutorial}
\textbf{Objectif.} Conduire l'adjudication synthétique à dix ans de
l'\cref{ex:m2:the-treasury-market:auction} à partir de son carnet de soumissions et calculer chaque
statistique que le marché lit à 13:01. \textbf{État final~:}
\cref{fig:m2:the-treasury-market:demand} et les chiffres de l'exemple.
\begin{tutsteps}
\item \textbf{Plafonnement et tri.} Les soumissions de chaque enchérisseur ne comptent que jusqu'à 35~\% de l'émission.
\omcode{code/firm/tsyauction/firm_tsyauction.py}{33}{42}{Reconnaissance des soumissions de chaque enchérisseur jusqu'au plafond.}\label{lst:m2:the-treasury-market:cap}
\item \textbf{Servir jusqu'au rendement de coupure}, avec le pourcentage au prorata arrondi par excès
comme le prévoit la réglementation.
\omcode{code/firm/tsyauction/firm_tsyauction.py}{45}{65}{L'adjudication à prix uniforme~: rendement de coupure, allocation au rendement de coupure, attributions.}\label{lst:m2:the-treasury-market:run}
\item \textbf{Lire le résultat.} \texttt{tsy\_auction\_demo.summary()} donne le
rendement de coupure de 4.198~\%, le \omterm{def:m2:the-treasury-market:tail}{tail d'adjudication} de 1.8 points de base par rapport au \omterm{def:m2:the-treasury-market:wi}{marché gris} (when-issued) de 4.180~\%,
le \omterm{def:m2:the-treasury-market:tail}{taux de couverture} (bid-to-cover) 1.82, l'allocation au rendement de coupure 95.02~\% et les parts
attribuées~; \texttt{coupon\_from\_high\_yield(0.04198)} donne 4.125.
\item \textbf{Tracer la courbe de demande} avec \texttt{fig\_tsy\_auction.py}.
\end{tutsteps}
\textbf{Pour aller plus loin.} Relever d'un point de base chaque soumission indirecte
et observer de combien se déplace le rendement de coupure (l'\cref{exo:m2:the-treasury-market:7})~; puis
confier à un enchérisseur 60~\% du carnet de soumissions et voir le plafond entrer en jeu.
\end{tutorial}

\section{Construction~: l'analyseur d'adjudications}\label{bld:m2:the-treasury-market:auction}

\begin{build}
\bfield{Rôle} Le desk de taux de la société miniature soumissionne aux adjudications et négocie
autour d'elles~; sa recherche lit les statistiques de chaque adjudication. Ce composant
exécute une adjudication à partir d'un carnet de soumissions et évalue un résultat par rapport
au \omterm{def:m2:the-treasury-market:wi}{marché gris} (when-issued).
\bfield{Interface} \texttt{Bid(bidder, yld, amount)}~;
\texttt{run\_auction(offering, noncompetitive, bids)} renvoyant \texttt{Result(stop, allotment\_at\_stop, awards, competitive\_tendered,
bid\_to\_cover)}~; \texttt{cap\_bids}~; \texttt{tail\_bp(stop, when\_issued)}~;
\texttt{cumulative\_demand}~; \texttt{coupon\_from\_high\_yield}.
\bfield{Règles} Non concurrentielles en premier~; plafond de 35~\% par enchérisseur~; rendement de coupure au
rendement où les soumissions cumulées atteignent le solde~; le pourcentage au prorata arrondi
par excès à 0.01~\%~; une adjudication non couverte est une erreur~; coupon arrondi à la baisse au
huitième, avec un minimum d'un huitième.
\bfield{Tests d'acceptation} \texttt{code/firm/tsyauction/tests/}~: chaque gagnant
paie le rendement de coupure~; les attributions non concurrentielles réduisent la capacité concurrentielle~; le
plafond~; un carnet non couvert~; le signe du tail~; la règle du coupon.
\bfield{Pour aller plus loin} Lire les résultats publiés des adjudications du Trésor et calculer
le tail par rapport à un flux when-issued~; ajouter les catégories de soumissionnaires que publie le
Trésor~; évaluer le nouveau titre avec la bibliothèque d'obligations
(\cref{bld:m2:government-bonds:bond}).
\end{build}

\begin{omsources}
\item US Treasury, \emph{How Auctions Work} et \emph{Auctions In Depth}
(TreasuryDirect)~; 31 CFR 356.20.
\item Federal Reserve Bank of New York, \emph{Primary Dealers}. \item US Department of the Treasury, Board of Governors of the Federal
Reserve System, Federal Reserve Bank of New York, SEC et CFTC, \emph{Joint
Staff Report: The U.S. Treasury Market on October 15, 2014}, juillet 2015. \item Financial Stability Board, \emph{Holistic Review of the March Market
Turmoil}, 17 novembre 2020.
\item Communiqué de presse de la Federal Reserve du 15 mars 2020~; Federal Reserve Bank of
New York, \emph{FAQs: Treasury Purchases}, 23 mars 2020~; H.4.1 via FRED.
\item SEC, communiqué de presse 2025-43~; M.\,T.\ Uyeda, propos prononcés lors de la
2026 U.S. Treasury Market Conference, 22 septembre 2026.
\end{omsources}

\section{Exercices}

\begin{exercise}[$\star$]\label{exo:m2:the-treasury-market:1}
Il est proposé 30~milliards de dollars~; les soumissions non concurrentielles totalisent 3~milliards de dollars. Quatre soumissionnaires
proposent 10~milliards de dollars à 4.100~\%, 8~milliards à 4.105~\%, 10~milliards à 4.110~\% et
6~milliards à 4.115~\%. Donner le rendement de coupure, l'allocation au taux de coupure et le
\omterm{def:m2:the-treasury-market:tail}{taux de couverture} (bid-to-cover).
\end{exercise}

\begin{exercise}[$\star$]\label{exo:m2:the-treasury-market:2}
Le rendement du \omterm{def:m2:the-treasury-market:wi}{marché gris} à 13~h est de 4.285~\%. Donner le tail si le rendement de coupure est
de 4.292~\%, puis s'il est de 4.281~\%. Quel résultat le marché jugerait-il fort~?
\end{exercise}

\begin{exercise}[$\star$]\label{exo:m2:the-treasury-market:3}
Donner le coupon d'une nouvelle obligation dont l'adjudication s'arrête à 4.198~\%, puis à 3.999~\%.
Cette nouvelle obligation est-elle émise au-dessus ou en dessous du pair~?
\end{exercise}

\begin{exercise}[$\star\star$]\label{exo:m2:the-treasury-market:4}
Un titre à dix ans \omterm{def:m2:the-treasury-market:otr}{on-the-run} rapporte 4.18~\% et le titre à dix ans précédent, trois mois
plus ancien, rapporte 4.21~\%. Donner trois raisons de cet écart et décrire l'opération qui
parie sur sa résorption. Qu'est-ce qui peut mal tourner~?
\end{exercise}

\begin{exercise}[$\star\star$]\label{exo:m2:the-treasury-market:5}
Dans l'\cref{ex:m2:the-treasury-market:auction}, quelle est la part au prorata d'un
\omterm{def:m2:the-treasury-market:pd}{spécialiste en valeurs du Trésor} dans l'émission~? Les spécialistes ont pris 38.3~\% de
l'adjudication. Que nous indique ce chiffre que le \omterm{def:m2:the-treasury-market:tail}{taux de couverture} ne dit pas au marché~?
\end{exercise}

\begin{exercise}[$\star\star$]\label{exo:m2:the-treasury-market:6}
Un investisseur unique souhaite 20~milliards de dollars d'une adjudication de 42~milliards. Quelle part
de sa soumission est retenue~? Comment peut-il malgré tout finir par détenir 20~milliards~?
\end{exercise}

\begin{exercise}[$\star\star\star$]\label{exo:m2:the-treasury-market:7}
\emph{Programmation.} Dans \texttt{tsy\_auction\_demo.bid\_book()}, augmenter d'un point de base le
rendement de chaque soumissionnaire indirect. Indiquer le nouveau rendement de coupure et le tail. Pourquoi
le rendement de coupure se déplace-t-il de moins, ou d'autant, que le décalage~?
\end{exercise}

\begin{exercise}[$\star\star\star$]\label{exo:m2:the-treasury-market:8}
\emph{Repérer l'erreur.} «~Dans une \omterm{def:m2:the-treasury-market:auction}{adjudication à prix uniforme}, il faut soumissionner un rendement
un peu supérieur à ce que l'on pense que vaut le titre, puisqu'on paie ce qu'on a soumis~». Corriger
cette affirmation. Qui a intérêt à réduire ses soumissions, et pourquoi~?
\end{exercise}

\section{Problème~: calcul du tail}

\begin{problem}[{Problème du week-end --- l'adjudication d'un spécialiste en valeurs du Trésor}]\label{pb:m2:the-treasury-market:1}
Le \omterm{def:m2:the-treasury-market:pd}{spécialiste en valeurs du Trésor} D07 participe à l'adjudication de l'\cref{ex:m2:the-treasury-market:auction}.
À 12~h~59, il a vendu à des clients 1~milliard de dollars du nouveau titre sur le \omterm{def:m2:the-treasury-market:wi}{marché gris} à 4.180~\%, et il
soumet 500~millions de dollars à 4.188~\%, 500~millions à 4.195~\% et 615
millions à 4.205~\%. Ses coûts aller-retour représentent un quart de point de base de
DV01. Le DV01 du nouveau titre au taux de coupure est de 807~dollars par million.

\medskip\noindent\textbf{Partie I --- L'adjudication.}
\begin{enumerate}
\item Quelle part de l'émission est disponible pour les soumissionnaires concurrentiels~?
\item Donner le rendement de coupure et le tail.
\item Donner le \omterm{def:m2:the-treasury-market:tail}{taux de couverture} et l'allocation au taux de coupure.
\item Donner les parts des soumissionnaires indirects, des soumissionnaires directs et des spécialistes.
\item Donner le coupon et le prix que paie chaque adjudicataire.
\end{enumerate}

\medskip\noindent\textbf{Partie II --- Le spécialiste.}
\begin{enumerate}[resume]
\item Quelle est la part au prorata de D07 dans l'émission~? Sa soumission respecte-t-elle
l'obligation~?
\item Lesquelles de ses soumissions sont retenues, et quel montant lui est attribué~?
\item Quel est le DV01 de son attribution~?
\item Elle est vendeuse de 1~milliard en \omterm{def:m2:the-treasury-market:wi}{marché gris} à 4.180~\% et acheteuse de 1~milliard au
taux de coupure. Quel est son P\&L après coûts~?
\item Quel aurait été son résultat si l'adjudication s'était arrêtée 0.2~point de base plus bas
(à 4.178~\%), à attribution identique~?
\end{enumerate}

\medskip\noindent\textbf{Partie III --- Lecture du résultat.}
\begin{enumerate}[resume]
\item À quel tail la position du spécialiste atteint-elle le seuil d'équilibre~?
\item Pourquoi le marché lit-il un grand tail comme une demande faible, et que fait le
titre sur le \omterm{def:m2:the-treasury-market:wi}{marché gris} dans la minute qui suit le résultat~?
\item Pourquoi le spécialiste a-t-il soumis une partie de sa part au-dessus du rendement du
\omterm{def:m2:the-treasury-market:wi}{marché gris}, au fond~?
\item Si les soumissionnaires indirects avaient tous soumissionné un point de base plus haut, où
l'adjudication se serait-elle arrêtée~?
\item Qu'aurait souhaité faire autrement le spécialiste en connaissant le tail à l'avance, et
pourquoi ne le peut-il pas~?
\end{enumerate}

\medskip\noindent\textbf{Partie IV --- Jugement.}
\begin{enumerate}[resume]
\item Pourquoi le Trésor utilise-t-il une \omterm{def:m2:the-treasury-market:auction}{adjudication à prix uniforme} plutôt qu'une
adjudication discriminatoire (pay-your-bid)~?
\item Pourquoi les clients qui ont acheté en \omterm{def:m2:the-treasury-market:wi}{marché gris} auprès de D07 ne soumissionnent-ils pas
simplement eux-mêmes à l'adjudication~?
\item Que coûte au spécialiste son obligation de soumission lors d'une adjudication faible,
et que lui rapporte-t-elle en échange~?
\item Énoncer le \emph{résultat nommé}~: le P\&L du spécialiste par point de base de tail sur
son attribution, son P\&L dans cette adjudication et le tail d'équilibre.
\item En une phrase~: qu'est-ce qu'un tail~?
\end{enumerate}
\end{problem}

\section{Questions d'entretien}

\begin{interviewq}[$\star$ \iqroles{trader, chercheur}]\label{iq:m2:the-treasury-market:1}
Qu'est-ce qu'un \omterm{def:m2:the-treasury-market:tail}{tail d'adjudication}, et que vous indique un tail élevé~?
\end{interviewq}

\begin{interviewq}[$\star$ \iqroles{trader, banque}]\label{iq:m2:the-treasury-market:2}
Pourquoi un Treasury \omterm{def:m2:the-treasury-market:otr}{on-the-run} rapporte-t-il en général moins qu'un \omterm{def:m2:the-treasury-market:otr}{off-the-run} de
maturité presque identique~?
\end{interviewq}

\begin{interviewq}[$\star\star$ \iqroles{trader, chercheur}]\label{iq:m2:the-treasury-market:3}
Comment une \omterm{def:m2:the-treasury-market:auction}{adjudication à prix uniforme} modifie-t-elle la stratégie d'un soumissionnaire par
rapport à une adjudication discriminatoire (pay-your-bid)~?
\end{interviewq}

\begin{interviewq}[$\star\star$ \iqroles{chercheur, développeur}]\label{iq:m2:the-treasury-market:4}
Que s'est-il passé sur le marché du Treasury le 15 octobre 2014, et que
chercheriez-vous dans les données du carnet d'ordres pour comprendre une journée comme celle-là~?
\end{interviewq}

\begin{interviewq}[$\star\star$ \iqroles{trader, banque}]\label{iq:m2:the-treasury-market:5}
En mars 2020, les Treasuries ont reculé alors que les actions
s'effondraient. Pourquoi l'actif le plus sûr a-t-il baissé~?
\end{interviewq}

\begin{interviewq}[$\star\star\star$ \iqroles{développeur, chercheur}]\label{iq:m2:the-treasury-market:6}
Concevoir le modèle de données et les tests d'un système qui ingère les résultats des
adjudications de Treasuries et calcule le tail de chaque adjudication en temps réel.
\end{interviewq}
